\subsection{Improper Integrals}

\begin{definition} \label{def:improper-riemann-integral} \hlt{Improper Riemann Integrals}\\
Let $a,b\in\overline{\R}$ satisfy $a<b$, and let $f:(a,b)\to\R$. Suppose $(\forall c,d\in\R)\,[a<c<d<b\Rightarrow f|_{[c,d]}\in\mathcal{R}[c,d]]$.\\
Fix $\xi\in(a,b)$. A one-sided integral below is defined whenever its limit exists in $\overline{\R}$.
\begin{enumerate}[label=(\roman*)]
\setlength{\itemsep}{0pt}
\item $\int_\xi^b f(x)\,dx\defeq\lim\limits_{d\uparrow b}\int_\xi^d f(x)\,dx$, and $\int_a^\xi f(x)\,dx\defeq\lim\limits_{c\downarrow a}\int_c^\xi f(x)\,dx$.
\item If their sum is defined in $\overline{\R}$, set $\int_a^b f(x)\,dx\defeq\int_a^\xi f(x)\,dx+\int_\xi^b f(x)\,dx$.\\
This value is independent of $\xi$ by Theorem \ref{thm:riemann-stieltjes-integral-properties}(iii) with the identity integrator.
\item The improper integral \hlt{converges} iff both one-sided integrals lie in $\R$; it \hlt{diverges} otherwise.
\end{enumerate}
\end{definition}

\noindent Figure \ref{fig:improper-integral-truncations} depicts the two right-endpoint truncations in Definition \ref{def:improper-riemann-integral}.\\
The left panel has $b\in\R$, while the right panel has $b=+\infty$.
\begin{illustrationfigure}[H]
  \centering
  \begin{tikzpicture}[
  x=0.78cm,
  y=0.72cm,
  >=Stealth,
  every node/.style={font=\scriptsize},
  axis/.style={->, black!55, line width=0.45pt},
  curve/.style={black!85, line width=0.9pt},
  area/.style={fill=gray!16, draw=none},
  guide/.style={black!30, dashed, line width=0.4pt}
]
  \begin{scope}
    \draw[axis] (-0.2,0) -- (5.05,0) node[right] {$x$};
    \draw[axis] (0,-0.15) -- (0,3.15);
    \path[area]
      (0,0) -- plot[domain=0:3.65, samples=100] (\x,{0.42+0.28*\x+0.07/((4.28-\x)*(4.28-\x))}) -- (3.65,0) -- cycle;
    \draw[curve, domain=0:4.02, samples=130]
      plot (\x,{0.42+0.28*\x+0.07/((4.28-\x)*(4.28-\x))});
    \draw[guide] (3.65,0) -- (3.65,2.35);
    \draw[guide] (4.28,0) -- (4.28,3.00);
    \node[below] at (0,0) {$a$};
    \node[below] at (3.65,0) {$d$};
    \node[below] at (4.28,0) {$b$};
    \node at (2.25,-0.55) {$d\to b^-$};
  \end{scope}

  \begin{scope}[xshift=6.4cm]
    \draw[axis] (-0.2,0) -- (5.7,0) node[right] {$x$};
    \draw[axis] (0,-0.15) -- (0,3.15);
    \path[area]
      (0,0) -- plot[domain=0:4.35, samples=120] (\x,{2.35/(1+\x)+0.18}) -- (4.35,0) -- cycle;
    \draw[curve, domain=0:5.25, samples=150]
      plot (\x,{2.35/(1+\x)+0.18});
    \draw[guide] (4.35,0) -- (4.35,0.62);
    \node[below] at (0,0) {$a$};
    \node[below] at (4.35,0) {$d$};
    \node[below] at (5.45,0) {$\infty$};
    \node at (2.80,-0.55) {$d\to\infty$};
  \end{scope}
\end{tikzpicture}

  \caption{Compact truncations approach a finite singular endpoint on the left and $+\infty$ on the right.}
  \label{fig:improper-integral-truncations}
\end{illustrationfigure}

\begin{theorem} \label{thm:nonnegative-improper-integral-supremum} \hlt{Nonnegative Improper Integrals}\\
Let $a\in\R$, $b\in\overline{\R}$, $a<b$, and let $f:[a,b)\to[0,+\infty)$ satisfy $(\forall d\in(a,b))\,f|_{[a,d]}\in\mathcal{R}[a,d]$.\\
Set $H(d)\defeq\int_a^d f(x)\,dx$ for $d\in(a,b)$ and $S\defeq\sup\limits_{\overline{\R}}H((a,b))$. Then:
\begin{enumerate}[label=(\roman*)]
\setlength{\itemsep}{0pt}
\item $\int_a^b f(x)\,dx=\lim\limits_{d\uparrow b}H(d)=S\in[0,+\infty]$.
\item The integral $\int_a^b f(x)\,dx$ converges iff $S\in\R$.
\end{enumerate}
\begin{proof}
\begin{enumerate}[label=(\roman*)]
\setlength{\itemsep}{0pt}
\item[$(i)$] Inequalities $0\leq H(d_1)\leq H(d_2)$ hold for $a<d_1<d_2<b$ by Theorem \ref{thm:riemann-stieltjes-integral-properties}(ii)--(iii) with $\alpha(t)=t$; $H$ increasing, $S\in[0,+\infty]$.\\
If $S\in\R$, then $(\forall\varepsilon>0)(\exists d_0\in(a,b))(\forall d\in[d_0,b))\,S-\varepsilon<H(d)\leq S$ by the definition of $S$ and monotonicity; hence $H(d)\to S$.\\
If $S=+\infty$, then $(\forall K\in\R)(\exists d_0\in(a,b))(\forall d\in[d_0,b))\,H(d)>K$ by the same properties; hence $H(d)\to+\infty$.\\
Fix $\xi\in(a,b)$. The limit $\lim\limits_{c\downarrow a}\int_c^\xi f(x)\,dx=H(\xi)$ holds by Theorems \ref{thm:integral-function-continuity-and-derivative}(i) and \ref{thm:riemann-stieltjes-integral-properties}(iii).\\
For $\xi<d<b$, $\int_\xi^d f(x)\,dx=H(d)-H(\xi)$ by Theorem \ref{thm:riemann-stieltjes-integral-properties}(iii) with $\alpha(t)=t$.\\
Thus $\int_\xi^b f(x)\,dx=S-H(\xi)$ by Definition \ref{def:improper-riemann-integral}(i),  Theorem \ref{thm:extended-limit-algebra}(ii); $\int_a^b f(x)\,dx=H(\xi)+[S-H(\xi)]=S$ by Definition \ref{def:improper-riemann-integral}(ii).
\item[$(ii)$] ($\Rightarrow$) If the integral converges, then $S\in\R$ by Part (i) and Definition \ref{def:improper-riemann-integral}(iii).
\item[$(ii)$] ($\Leftarrow$) If $S\in\R$, then the integral converges by Part (i) and Definition \ref{def:improper-riemann-integral}(iii).
\end{enumerate}
\hlt{Structural Commentary}: Monotonicity fixes the extended limit; additivity transfers it to the improper integral.\\
\hlt{Key Insight}: Bounded truncation areas are exactly finite improper mass.
\end{proof}
\end{theorem}

\begin{perspective} \label{per:cauchy-principal-value-not-improper-integral} \hlt{Principal Value Is Different}\\
Definition \ref{def:improper-riemann-integral} requires separate one-sided limits on $(-\infty,+\infty)$. For $A>0$, $\int_{-A}^{A}\sin x\,dx=0$, so the Cauchy principal value is $0$.\\
However, $\int_0^d\sin x\,dx=1-\cos d$ has no limit as $d\to+\infty$, so the improper integral diverges.
\end{perspective}

\begin{definition} \label{def:improper-riemann-stieltjes-real-line} \hlt{Improper Riemann--Stieltjes Integrals on $\R$}\\
Let $\alpha:\R\to\R$ be increasing. Let $f:\R\to\R$ satisfy $(\forall a,b\in\R)\,[a<b\Rightarrow f|_{[a,b]}\in\mathcal{R}(\alpha|_{[a,b]})]$.
\begin{enumerate}[label=(\roman*)]
\setlength{\itemsep}{0pt}
\item Set $\alpha(-\infty)\defeq\inf\limits_{\overline{\R}}\alpha(\R)$ and $\alpha(+\infty)\defeq\sup\limits_{\overline{\R}}\alpha(\R)$.
\item When the limits exist in $\overline{\R}$, set $I_-(f,\alpha)\defeq\lim\limits_{a\to-\infty}\int_a^0 f\,d\alpha$. Under the same condition, set $I_+(f,\alpha)\defeq\lim\limits_{b\to+\infty}\int_0^b f\,d\alpha$.
\item If $I_-(f,\alpha)+I_+(f,\alpha)$ is defined in $\overline{\R}$, set $\int_{-\infty}^{+\infty}f\,d\alpha\defeq I_-(f,\alpha)+I_+(f,\alpha)$.
\item The function $f$ is \hlt{$\alpha$-integrable on $\R$} iff $I_-(f,\alpha),I_+(f,\alpha)\in\R$.
\end{enumerate}
\end{definition}

\begin{theorem} \label{thm:improper-riemann-stieltjes-nonnegative-bounded-mass} \hlt{Nonnegative and Bounded-Mass Stieltjes Integrals}\\
Let $\alpha:\R\to\R$ be increasing. Let $f:\R\to\R$ satisfy $(\forall a,b\in\R)\,[a<b\Rightarrow f|_{[a,b]}\in\mathcal{R}(\alpha|_{[a,b]})]$.\\
Set $S_+\defeq\sup\limits_{\overline{\R}}\{\int_0^b f\,d\alpha:b>0\}$ and $S_-\defeq\sup\limits_{\overline{\R}}\{\int_a^0 f\,d\alpha:a<0\}$.\\
For Parts (i)--(iii), assume $(\forall x\in\R)\,f(x)\geq0$; Part (iv) has no sign hypothesis. Then:
\begin{enumerate}[label=(\roman*)]
\setlength{\itemsep}{0pt}
\item $I_-(f,\alpha)=S_-$ and $I_+(f,\alpha)=S_+$ belong to $[0,+\infty]$.
\item $\int_{-\infty}^{+\infty}f\,d\alpha=S_-+S_+\in[0,+\infty]$.
\item The function $f$ is $\alpha$-integrable on $\R$ iff $S_-,S_+\in\R$.
\item If $\alpha(-\infty),\alpha(+\infty)\in\R$ and $f$ is bounded, then $f$ is $\alpha$-integrable on $\R$.
\end{enumerate}
\begin{proof}
\begin{enumerate}[label=(\roman*)]
\setlength{\itemsep}{0pt}
\item[$(i)$] Set $H_+(t)\defeq\int_0^t f\,d\alpha$ and $H_-(t)\defeq\int_{-t}^0 f\,d\alpha$ for $t>0$. For $0<s<t$, $0\leq H_\pm(s)\leq H_\pm(t)$ by Theorem \ref{thm:riemann-stieltjes-integral-properties}(ii)--(iii).\\
Fix $\sigma\in\{-,+\}$. If $S_\sigma\in\R$, then $(\forall\varepsilon>0)(\exists t_0>0)(\forall t\geq t_0)\,S_\sigma-\varepsilon<H_\sigma(t)\leq S_\sigma$ by the supremum property and monotonicity.\\
If $S_\sigma=+\infty$, then $(\forall K\in\R)(\exists t_0>0)(\forall t\geq t_0)\,H_\sigma(t)>K$ by the same properties.\\
Thus $H_\sigma(t)\to S_\sigma$, and $I_\pm(f,\alpha)=S_\pm$ by Definition \ref{def:improper-riemann-stieltjes-real-line}(ii), with $a=-t$ for $\sigma=-$.
\item[$(ii)$] The identity $\int_{-\infty}^{+\infty}f\,d\alpha=S_-+S_+$ holds by Part (i) and Definition \ref{def:improper-riemann-stieltjes-real-line}(iii).
\item[$(iii)$] ($\Rightarrow$) If $f$ is $\alpha$-integrable, then $S_-,S_+\in\R$ by Part (i) and Definition \ref{def:improper-riemann-stieltjes-real-line}(iv).
\item[$(iii)$] ($\Leftarrow$) If $S_-,S_+\in\R$, then $f$ is $\alpha$-integrable by Part (i) and Definition \ref{def:improper-riemann-stieltjes-real-line}(iv).
\item[$(iv)$] Choose $M>0$ with $(\forall x\in\R)\,\abs{f(x)}\leq M$, and set $g_\pm\defeq M\pm f$.\\
Each $g_\pm$ is locally Riemann--Stieltjes integrable by Definition \ref{def:riemann-stieltjes-integral} and Theorem \ref{thm:riemann-stieltjes-integral-properties}(i).\\
The bounds $(\forall x\in\R)\,0\leq g_\pm(x)\leq2M$ hold by the choice of $M$.\\
For $g\in\{g_-,g_+\}$, $0\leq I_+(g,\alpha)\leq2M[\alpha(+\infty)-\alpha(0)]\in\R$ by Part (i) and Theorem \ref{thm:riemann-stieltjes-integral-properties}(iv).\\
For $g\in\{g_-,g_+\}$, $0\leq I_-(g,\alpha)\leq2M[\alpha(0)-\alpha(-\infty)]\in\R$ by the same part and theorem.\\
Thus $I_\pm(f,\alpha)=\frac{1}{2}[I_\pm(g_+,\alpha)-I_\pm(g_-,\alpha)]\in\R$ by Theorems \ref{thm:riemann-stieltjes-integral-properties}(i), \ref{thm:extended-limit-algebra}(ii)--(iii); $f$ $\alpha$-integrable by Definition \ref{def:improper-riemann-stieltjes-real-line}(iv).
\end{enumerate}
\hlt{Structural Commentary}: Suprema identify nonnegative limits; the shifts $M\pm f$ reduce bounded integrands to them.\\
\hlt{Key Insight}: Finite integrator mass bounds both positive shifts uniformly.
\end{proof}
\end{theorem}
